\documentclass[10pt]{article}
\usepackage[T1]{fontenc}
\usepackage[utf8]{inputenc}
\usepackage[french]{babel}
\usepackage{amsmath,amssymb,mathrsfs,mathpazo}
\usepackage{geometry}
\geometry{a4paper,margin=2cm}
\usepackage{graphicx,array,tabularx,tabularray,colortbl,multirow,multicol,booktabs,enumitem,pifont,ulem,fancybox,diagbox,cancel}
\usepackage{pstricks,pstricks-add,pst-plot,pst-tree,pst-node,pst-text,pst-func,pst-math,pst-bspline,pst-3dplot,pst-eucl,pst-solides3d}
\let\pscalculate\undefined
\usepackage{pst-bezier}
\usepackage{tikz,pgfplots,tkz-tab}
\usetikzlibrary{arrows,arrows.meta,patterns,calc,positioning,trees,decorations.markings,shapes.geometric,shapes.misc,shapes.arrows,intersections,decorations.pathmorphing,decorations.pathreplacing,backgrounds,shadows,babel}
\pgfplotsset{compat=1.18}
\usepgfplotslibrary{fillbetween}
\usepackage[np]{numprint}
\usepackage{listings,makecell,pas-tableur,twoopt}
\let\barre\undefined
\usepackage[pstricks]{tablvar}
\definecolor{darkgreen}{RGB}{0,100,0}
\providecommand{\red}{\color{red}}
\providecommand{\blue}{\color{blue}}
\providecommand{\green}{\color{green}}
\providecommand{\black}{\color{black}}
\providecommand{\white}{\color{white}}
\providecommand{\gray}{\color{gray}}
\providecommand{\cyan}{\color{cyan}}
\providecommand{\magenta}{\color{magenta}}
\providecommand{\yellow}{\color{yellow}}
\providecommand{\violet}{\color{violet}}
\providecommand{\orange}{\color{orange}}
\providecommand{\euro}{\texteuro}
\providecommand{\EUR}{\texteuro}
\providecommand{\ieme}{\textsuperscript{e}}
\providecommand{\degres}{\ensuremath{^\circ}}
\providecommand{\degree}{\ensuremath{^\circ}}
\providecommand{\vv}[1]{\overrightarrow{#1}}
\providecommand{\Coord}[1]{#1}
\providecommand{\decosix}{}
\providecommand{\starredbullet}{\ensuremath{\bullet}}
\pagestyle{empty}
\setlength{\parindent}{0pt}
\usepackage{adjustbox}\DeclareUnicodeCharacter{2236}{:}\DeclareUnicodeCharacter{2190}{\ensuremath{\leftarrow}}
\providecommand{\R}{\mathbb{R}}
\providecommand{\N}{\mathbb{N}}
\providecommand{\Z}{\mathbb{Z}}
\providecommand{\Q}{\mathbb{Q}}
\providecommand{\D}{\mathbb{D}}
\providecommand{\C}{\mathbb{C}}
\providecommand{\e}{{\mathrm{e}}}
\providecommand{\vect}[1]{\overrightarrow{#1}}
\providecommand{\vectt}[1]{\overrightarrow{\mathrm{#1}}}
\providecommand{\vvt}[1]{\overrightarrow{\mathrm{#1}}}
\providecommand{\np}[2][]{{#2~#1}}
\providecommand{\vv}[1]{{\overrightarrow{#1}}}
\renewcommand{\d}{{\mathrm{d}}}
\providecommand{\degres}{{^\circ}}
\providecommand{\degre}{{^\circ}}
\providecommand{\blue}{\color{blue}}
\providecommand{\red}{\color{red}}
\providecommand{\enskip}{\quad}
\providecommand{\llbracket}{\lbrack\!\lbrack}
\providecommand{\rrbracket}{\rbrack\!\rbrack}
\providecommand{\euro}{{\text{€}}}
\providecommand{\no}{{\text{n°}}}
\providecommand{\interval}[2]{{[#1;#2]}}
\providecommand{\intervFO}[2]{{[#1;#2[}}
\providecommand{\dotfill}{\ldots}

\providecommand{\ds}{\displaystyle}
\providecommand{\textding}[1]{#1}
\providecommand{\Oij}{\(\mathrm{O};\vec{\imath},\vec{\jmath}\)}
\providecommand{\Oijk}{\(\mathrm{O};\vec{\imath},\vec{\jmath},\vec{k}\)}
\providecommand{\Ouv}{\(\mathrm{O};\vec{u},\vec{v}\)}
\providecommand{\intervFF}[2]{\ensuremath{[#1;#2]}}
\renewcommand{\euro}{\ifmmode\text{\texteuro}\else\texteuro\fi}
\DeclareUnicodeCharacter{00B0}{\ifmmode^{\circ}\else\textdegree\fi}
\DeclareUnicodeCharacter{20AC}{\ifmmode\text{\texteuro}\else\texteuro\fi}
% Énoncés : Éducation nationale. Transcription et corrigés : APMEP, Denis Vergès et contributeurs. Découpage et classement : CoopMaths, Sébastien Lozano, Éric Elter et contributeurs. Adaptation de la présentation : Maths Explorer.
% CC BY-SA : https://coopmaths.fr/www/about/
% https://coopmaths.fr/alea/static/e3c/2020/tex/e3c_2020_00_sujet32_1_cor.tex
\begin{document}
\noindent\textit{APMEP / CoopMaths. CC BY-SA. Présentation adaptée par Maths Explorer.}\par\medskip
\subsection*{Question 1}

Le produit est nul si l'un des facteurs est nul :
\[
x - 1 = 0 \text{ si } x = 1,
\]
\[
x^2 + x + 1 = 0. \text{ On a } \Delta = 1 - 4 = -3 < 0,
\]
ce trinôme n'a donc pas de racines.

L'équation a donc une seule solution : \(1\).

\subsection*{Question 2}

Comme, quel que soit le réel \(x\), \(\e^x > 0\), on a \(\e^x + 1 > 1 > 0\), \(f(x)\) ne peut s'annuler que si : \[
7x - 23 = 0 \iff x = \dfrac{23}{7}.
\]

\subsection*{Question 3}

le cercle de centre \(A(-4\,;\,2)\) et de rayon \(r = \sqrt{2}\) a pour équation :
\begin{align*}
&M(x\,;\,y) \in \mathcal{C}(A\,;\,R = \sqrt{2}) \\
\iff &AM^2 = (\sqrt{2})^2 \\
\iff &(x + 4)^2 + (y - 2)^2 = 2.
\end{align*}

\subsection*{Question 4}
 
\(\vec{u}\) et \(\vec{v}\) sont orthogonaux si et seulement si leur produit scalaire est nul, soit si :
\begin{align*}
&\vec{u} \cdot \vec{v} = 0 \\
\iff &m(m + 1) - 2 = 0 \\
\iff &m^2 + m - 2 = 0.
\end{align*}
Pour le trinôme, \(\Delta = 1 + 8 = 9 = 3^2 > 0\), il a donc deux racines :
\[
m_1 = \dfrac{-1 + 3}{2} = 1 \quad \text{et} \quad m_2 = \dfrac{-1 - 3}{2} = -2.
\]
Réponse : \(\textbf{b.}\)

\subsection*{Question 5}

Une équation cartésienne de la droite \(\mathcal{D}\) passant par le point \(A(-2\,;\,5)\) et admettant pour vecteur normal \(\vec{n} \begin{pmatrix} -1 \\ 3 \end{pmatrix}\) est :
\begin{align*}
&M(x\,;\,y) \in \mathcal{D} \\
\iff &\overrightarrow{AM} \cdot \vec{n} = 0 \\
\iff &-1(x - (-2)) + 3(y - 5) = 0 \\
\iff &-x - 2 + 3y - 15 = 0 \\
\iff &-x + 3y - 17 = 0 \\
\iff &x - 3y + 17 = 0.
\end{align*}

\end{document}
